\section{Primary descent from frames}
\label{sec:frame-primary-descent}
We record explicitly the descent used in the global primary
statements. Its hypotheses are verified on Zariski frame charts;
no claim about arbitrary regular-fibre morphisms is needed.

\begin{lemma}[Frame descent of primary decompositions]
\label{lem:frame-primary-descent}
Let $X$ be a noetherian scheme, let $\mathcal K$ be a vector bundle,
and let $f:F=\operatorname{Fr}(\mathcal K)\to X$ be its frame
torsor. Suppose $\mathcal I,\mathcal Q_1,\ldots,\mathcal Q_s$
are coherent ideals of $\mathcal O_X$ and
$Y_1,\ldots,Y_s\subset X$ are distinct integral closed subschemes.
Assume that on $F$ the pulled-back ideals give a primary decomposition
\[
 f^*\mathcal I=\bigcap_i f^*\mathcal Q_i,
 \qquad \sqrt{f^*\mathcal Q_i}=\mathcal I_{f^{-1}Y_i},
\]
with each nonempty local component primary. Then
$\mathcal I=\bigcap_i\mathcal Q_i$ on $X$, each $\mathcal Q_i$
is primary along $Y_i$, and the associated points of
$\mathcal O_X/\mathcal I$ are the generic points of the components
in its irredundant decomposition. The intersection on $X$ is
irredundant if and only if the intersection on $F$ is irredundant.

The hypotheses can be checked by polynomial extension and localization
of affine primary decompositions, provided every claimed associated
stratum meets the frame open. For an invariant symbolic-power ideal,
its intrinsic description on $X$ is its symbolic power at the generic
point of the indicated integral support.
\end{lemma}
\begin{proof}
On an affine trivializing open $U=\operatorname{Spec}A$, one has
\[
 f^{-1}U=\operatorname{Spec}A[z_{ij},\det(z)^{-1}].
\]
This ring is faithfully flat over $A$, and the inclusion has the
retraction $z_{ij}\mapsto\delta_{ij}$. Flatness commutes with a
finite intersection of ideals, as follows by tensoring the kernel
sequence defining the intersection. Faithful flatness detects ideal
equality, so the displayed intersection contracts to the asserted
one on $U$.

If $ab\in\mathcal Q_i(U)$ and $b\notin\mathcal Q_i(U)$,
faithfulness implies that $b$ remains outside its extended ideal.
Primarity upstairs therefore implies that a power of $a$ lies in
the extended ideal, and contraction puts that power in
$\mathcal Q_i(U)$. The radical contracts to the ideal of $Y_i$.
When a trivializing open requires several affine subcharts of $F$,
the same argument is local: an element nonzero in the quotient
remains nonzero on some such chart; the prime support specifies its
unique generic point. This proves the local primary assertion.
Conversely polynomial extension of a primary ideal is primary, and
localization preserves primarity unless the ideal becomes the unit
ideal. One may verify the polynomial assertion by the usual
zero-divisor criterion for polynomials over a primary ring, or by
contracting from the localization at the extended prime.

If a component can be omitted upstairs, faithful flatness contracts
the resulting equality. If it can be omitted downstairs, flatness
extends the equality. This proves the irredundancy equivalence.
On an affine chart the associated primes of an irredundant primary
decomposition with distinct radicals are exactly those radicals;
applying this statement and gluing gives the associated points on
$X$. In particular the assertion is not deduced solely from a
statement about fibre regularity.

For symbolic powers, polynomial extension followed by localization
at the extended prime commutes with taking the contraction of a
power of the maximal ideal. Equivalently apply the defining
membership test at the generic point on each trivialization.
Thus the ideals described this way agree on overlaps with the
intrinsic symbolic-power sheaves. All assertions are compatible
with further Zariski localization.
\end{proof}

\begin{corollary}[The global ideals in Hilbert function $(1,2,2)$]
\label{cor:loewy-primary-descent}
In Theorem~\ref{thm:loewy-primary}, let $\mathcal D=\mathcal I_\Delta$
and let $\mathcal P,\mathcal Q$ denote the ideals of the factor
surfaces $Y_L$ (one surface in the square case, two in the distinct
factor case). In the latter case put
\[
 \mathcal N=\mathcal P+\mathcal Q=\mathcal I_{\{K=S_2\}},\qquad
 \mathcal J=(\mathcal I_D:\mathcal D^2),
\]
\[
 \mathcal Q_0=\mathcal J+\mathcal P^2+\mathcal Q^2,
 \qquad
 \mathcal Q_*=\mathcal D^2\mathcal Q_0+\mathcal N^7.
\]
These are global coherent ideals on $\Gr(2,V\oplus S_2)$.
The square decomposition is
$\mathcal D^2\cap\mathcal P^4$; the distinct-factor decomposition
is $\mathcal D^2\cap\mathcal P^3\cap\mathcal Q^3\cap\mathcal Q_*$.
Both are irredundant globally and their associated points are
exactly those stated in Theorem~\ref{thm:loewy-primary}.
\end{corollary}
\begin{proof}
The surface ideal is the zero ideal of the bundle map
$K\to V/L$, and the point ideal is the zero ideal of $K\to V$.
Thus the ideals $\mathcal P,\mathcal Q,\mathcal N$ are intrinsic;
on frames they are $(a,c),(b,d),(a,b,c,d)$ after choosing the
factor lines as axes. The determinant ideal is intrinsic as well.
The ideal quotient by its square commutes with flat pullback:
locally that square is generated by a nonzerodivisor, and the
quotient is the kernel of multiplication by that generator on
$\mathcal O_X/\mathcal I_D$. This identifies its pullback with
the ideal $J$ used in the affine calculation. The stated sums and
products therefore pull back to $Q_0,Q_*$ exactly.

The affine primary equalities are proved in
Theorem~\ref{thm:loewy-primary}. Polynomial extension by the lower
frame block and restriction to the full-frame open preserve them.
Every associated stratum survives this restriction: the lower block
can be the identity matrix, independently of the upper block.
The strata $\Delta,Y_L,\{K=S_2\}$ are integral, and their inverse
images under a frame torsor are integral. Apply
Lemma~\ref{lem:frame-primary-descent} to obtain equality, primarity,
irredundancy and the exact associated-point set on the Grassmannian.
Permuting the two factor lines permutes the two surface ideals and
leaves the unordered decomposition unchanged. Thus no labelled
choice is needed for the asserted geometric collection.
\end{proof}
